\chapter{Motivation}
Im September 2025 sah man plötzlich im ganzen Internet Nachrichten zur Child Sexual Abuse Regulation (CSAR) der EU oder einfach EU Chat Control 2.0. In diesem Gesetzentwurf geht es darum, dass in den EU Staaten jede einzelne versendete Private nachricht und Datei von einem programm geprüft wird auf gefährliche Inhalte. Selbstverstädlich machte dies vielen Leuten angst um ihre Privatsphäre da es Ende-zu-Ende-Verschlüsselungen umging und eine permanente Überwachung aller EU Bürger erlaubte. Die andere Sorge war das unabhängige Staaten danach auch eigene Gesetze zur überwachung von online Kommunikation erlassen würden. Tim und ich haben auch über dieses Thema disskutiert und kamen zum schluss das wir zusammen eine eigene Platform zum chatten programmieren wollen auch wenn wir nicht direkt betroffen währen von der CSAR. Am nächsten Morgen haben wir ganz Tutti und Ricardo durckkämmt auf der suche nach einem alten PC welchen wir zu einem Linux (Arch btw) Homeserver umbauen können. Eine Woche später und 50 CHF ärmer waren wir in Rüti ZH um einen Lenovo Office-PC abzuholen. 
\\\\
Anfangs Januar hatten wir noch keine Zeile Code geschreiben. Wir hatten jedoch einen Plan wie wir zwei Fliegen mit einer Klappe schlagen konnten, denn Tim und ich hatten beide noch keine Ahnung zu welchen Thema wir unsere Maturitätsarbeit wollten. Unserem Plan folgend fragten Thomas Graf ob es möglich ist die Maturitätsarbeit zu zweit zu schreiben. Als es ja sagte war es für uns beide klar: dieses Projekt wird zu unserer Maturitätsarbeit. Zusammen schreiben wir in Python ein möglichst sicheres Chat Programm welches wir dann auf unserem Server hosten können. Somit haben wir sicher die motivation das Projekt abzuschliessen oder zumindest den Druck. Das einzige das uns noch fehlte war einen Schwerpunkt für die Arbeit. Auf einer langen Zugfahrt nach Interlaken sammelten wir viele Ideen die wir sofort wider verworfen, der grund dafür war meistens das sie zu viel Mathematik involvierten. Im endeffekt kamen wir zum schluss das wir uns damit beschäftigen wollen wie man die Menge an sammelbaren Metadaten minimalisieren oder sogar komplett eliminieren kann. Da dies ein sehr aktuelles Thema ist und Metadaten zu sammeln die einfachste art ist private Informationen aus chatplatformen zu gewinnen.